% ----------------------------------------------------------
% Fundamentação Teórica
% ----------------------------------------------------------
\chapter{Fundamentação Teórica} \label{cha:fundamentacao}

\section{Código Penal (CP)}
O Código Penal (CP) brasileiro, promulgado pelo Decreto-Lei nº 2.848/1940, estabelece as normas e princípios do Direito Penal, bem como define os tipos penais e as respectivas penas aplicáveis aos indivíduos que cometem crimes. O CP prevê três regimes básicos de cumprimento de pena: fechado, semiaberto e aberto, conforme o art.~33 do Código Penal. A definição do regime inicial é determinada pela quantidade de pena aplicada, das condições judiciais e pela reincidência do réu\cite{codigoPenal}.

A reincidência ocorre quando o apenado comete um novo crime, logo após a condenação anterior, já transitada em julgado – momento em que a decisão judicial se torna definitiva, não cabendo mais recursos. Assim, a sentença passa a ser definitiva e não pode mais ser alterada. A reincidência pode aumentar a pena vigente e tem duração de cinco anos, contados a partir do dia em que a pena do crime anterior foi executada ou extinta.\cite{codigoPenal}

Com a determinação da coisa julgada, os envolvidos devem cumprir com as obrigações impostas, sob as consequências legais da pena definida. Isso inclui as normas a serem executadas e a concessão dos benefícios previstos, à medida que a pena é realizada corretamente.\cite{codigoPenal}

Um cálculo de execução penal realizado de forma clara e objetiva, assegura que o condenado cumpra exatamente a sanção imposta na sentença, sendo coerente com os limites e critérios estabelecidos pela lei. Além de buscar a distinção da pena, ou seja, assegurar que o cumprimento da sentença seja de acordo com a situação específica de cada apenado.\cite{codigoPenal}

Quaisquer erros no cálculo podem atrasar indevidamente a progressão para um regime mais rigoroso, limitar o direito a saídas temporárias ou ao livramento condicional, o que configura excesso na execução, proibido pela lei.\cite{codigoPenal}

\begin{citacao}
Art. 33 - A pena de reclusão deve ser cumprida em regime fechado, semiaberto ou aberto. A de detenção, em regime semi-aberto, ou aberto, salvo necessidade de transferência a regime fechado. (Redação dada pela Lei nº 7.209, de 11.7.1984)

§ 1º - Considera-se: (Redação dada pela Lei nº 7.209, de 11.7.1984)
a) regime fechado a execução da pena em estabelecimento de segurança máxima ou média;
b) regime semi-aberto a execução da pena em colônia agrícola, industrial ou estabelecimento similar;
c) regime aberto à execução da pena em casa de albergado ou estabelecimento adequado.

§ 2º - As penas privativas de liberdade deverão ser executadas em forma progressiva, segundo o mérito do condenado, observados os seguintes critérios e ressalvadas as hipóteses de transferência a regime mais rigoroso:
a) o condenado à pena superior a 8 (oito) anos deverá começar a cumpri-la em regime fechado;
b) o condenado não reincidente, cuja pena seja superior a 4 (quatro) anos e não exceda a 8 (oito), poderá, desde o princípio, cumpri-la em regime semi-aberto;
c) o condenado não reincidente, cuja pena seja igual ou inferior a 4 (quatro) anos, poderá, desde o início, cumpri-la em regime aberto.

§ 3º - A determinação do regime inicial de cumprimento da pena far-se-á com observância dos critérios previstos no art. 59 deste Código.

§ 4º - O condenado por crime contra a administração pública terá a progressão de regime condicionada à reparação do dano que causou ou à devolução do produto do ilícito praticado, com os acréscimos legais.
(Código Penal, art. 33, \cite{codigoPenal})
\end{citacao}

Dessa forma, o juiz sentenciante fixa a pena base de acordo com o crime cometido. Posteriormente, o juiz aumenta ou diminui a pena, segundo o agravante (art. 61 do CP) e/ou atenuante (art. 65 do CP) do acusado, que modifica a pena base. Na fase seguinte, são analisadas as causas de aumento ou diminuição, que são fatores específicos de cada crime. Assim, analisadas às fases, é chegada a pena final com o tempo de pena e qual o regime a ser imposto inicialmente. Depois, quem passa a ficar responsável por acompanhar o cumprimento da pena é o juízo da execução penal.\cite{codigoPenal}

\section{Lei de Execução Penal(LEP)}
A Lei de Execução Penal (LEP), lei brasileira nº 7.210/1984, tem como finalidade garantir tanto o cumprimento da pena quanto os direitos do condenado. Como mencionado anteriormente neste capítulo, logo após o fim dos recursos para a condenação penal, o processo entra na fase de execução da pena, momento em que é regido pela LEP.\cite{lep}

\begin{citacao}
Art. 1º A execução penal tem por objetivo efetivar as disposições de sentença ou decisão criminal e proporcionar condições para a harmônica integração social do condenado e do internado.
(Lei de Execução Penal, art. 1, \cite{lep})
\end{citacao}

Basicamente trata-se das regras de como as penas privativas de liberdade (prisão) são executadas, como funcionam as penas restritivas de direitos, como deve ser a assistência ao preso, quais são os direitos e deveres das pessoas condenadas, como funciona a progressão de regime, como funciona benefícios como livramento condicional, indulto, remição de pena, etc.\cite{lep}

Após conclusão do processo que antecede a abertura da execução penal, em que o prazo de recursos é encerrado (trânsito em julgado), o caso é destinado ao Juiz da Execução Penal, que passa a acompanhar o cumprimento da pena. Ele recebe a guia de recolhimento (ou guia de execução), documento que contém todas as informações a respeito do condenado, bem como a pena imposta.\cite{lep}

Desse modo, a destinação do caso ao Juiz da Execução Penal representa a transição da fase de julgamento para a fase de garantia da efetivação da sentença penal, assegurando que a pena imposta seja realizada de forma justa e legal.\cite{lep}

\section{Processo de Execução Criminal(PEC)}
Depois que o juiz criminal determinar o número de anos, tudo o que modifica na prática o tempo que o condenado fica preso é tratado na Execução Penal. Para elaboração do cálculo de pena, obtém-se a data-base da pena, o regime inicial, os prazos de progressão de regime, livramento condicional, eventual indulto, extinção da pena, detração e unificação de pena, caso haja condenações anteriores ou posteriores. Assim, após ter esse cálculo por completo é que de fato a execução da pena do apenado começa a contar.\cite{lep}

O primeiro elemento de início de execução é o regime. A progressão de regime, que é a passagem de um regime mais rigoroso(fechado) para um menos rigoroso, ocorre de acordo com os requisitos objetivo e subjetivo (art. 112 da LEP). O requisito subjetivo depende da boa conduta e comportamento do apenado no estabelecimento prisional em que está encarcerado. Os requisitos objetivos estão relacionados ao lapso temporal, ou seja, ao intervalo de tempo que o apenado deve cumprir da pena total para ter direito à progressão de regime. Dessa forma, para avançar de regime é necessário cumprir um percentual do tempo total da pena, que varia conforme o crime cometido e a condição de reincidência.\cite{lep}

\begin{citacao}
Art. 112. A pena privativa de liberdade será executada em forma progressiva com a transferência para regime menos rigoroso, a ser determinada pelo juiz, quando o preso tiver cumprido ao menos:

I - 16\% (dezesseis por cento) da pena, se o apenado for primário e o crime tiver sido cometido sem violência à pessoa ou grave ameaça;

II - 20\% (vinte por cento) da pena, se o apenado for reincidente em crime cometido sem violência à pessoa ou grave ameaça;

III - 25\% (vinte e cinco por cento) da pena, se o apenado for primário e o crime tiver sido cometido com violência à pessoa ou grave ameaça;

IV - 30\% (trinta por cento) da pena, se o apenado for reincidente em crime cometido com violência à pessoa ou grave ameaça;

V - 40\% (quarenta por cento) da pena, se o apenado for condenado pela prática de crime hediondo ou equiparado, se for primário;

VI - 50\% (cinquenta por cento) da pena, se o apenado for:

a) condenado pela prática de crime hediondo ou equiparado, com resultado morte, se for primário, vedado o livramento condicional;

b) condenado por exercer o comando, individual ou coletivo, de organização criminosa estruturada para a prática de crime hediondo ou equiparado; ou

c) condenado pela prática do crime de constituição de milícia privada;

VI-A – 55\% (cinquenta e cinco por cento) da pena, se o apenado for condenado pela prática de feminicídio, se for primário, vedado o livramento condicional;

VII - 60\% (sessenta por cento) da pena, se o apenado for reincidente na prática de crime hediondo ou equiparado;

VIII - 70\% (setenta por cento) da pena, se o apenado for reincidente em crime hediondo ou equiparado com resultado morte, vedado o livramento condicional.

§ 1º Em todos os casos, o apenado somente terá direito à progressão de regime se ostentar boa conduta carcerária, comprovada pelo diretor do estabelecimento, e pelos resultados do exame criminológico, respeitadas as normas que vedam a progressão.

\end{citacao}

Posteriormente, caso haja detração — aproveitamento do tempo em que o apenado ficou detido provisoriamente antes da condenação definitiva, esse período será descontado no cálculo da pena-base. De acordo com o artigo 111° do Código Penal, caso ocorra reincidência na soma ou na unificação das penas, deve-se considerar o benefício existente, bem como a detração para definir o regime a ser cumprido.\cite{lep}

Durante o cumprimento da pena, o apenado pode usar do benefício de remição de pena, para que possa diminuir o tempo de cumprimento de pena, por meio de trabalho, estudo ou leitura.  Para remição por trabalho, o condenado no exercício do regime fechado ou semiaberto tem direito a 1 dia remido a cada 3 dias trabalhados. Para remição por estudo, 1 dia de pena a cada 12 horas de estudo de frequência escolar distribuídas em pelo menos 3 dias. Para remição por leitura, reduz-se 4 dias de pena a cada obra lida, sendo permitida no máximo 12 obras por ano (a depender das regras  normativas do tribunal). Dessa forma, o juiz da vara de execução criminal tem a responsabilidade de acompanhar de acordo com a administração penitenciária acerca do comportamento e andamento do apenado dentro do sistema prisional.\cite{lep}

\begin{citacao}
Art. 126. O condenado que cumpre a pena em regime fechado ou semiaberto poderá remir, por trabalho ou por estudo, parte do tempo de execução da pena.

§ 1º A contagem de tempo referida no caput será feita à razão de:

I - 1 (um) dia de pena a cada 12 (doze) horas de frequência escolar - atividade de ensino fundamental, médio, inclusive profissionalizante, ou superior, ou ainda de requalificação profissional - divididas, no mínimo, em 3 (três) dias;

II - 1 (um) dia de pena a cada 3 (três) dias de trabalho.

§ 2º As atividades de estudo a que se refere o § 1º deste artigo poderão ser desenvolvidas de forma presencial ou por metodologia de ensino a distância e deverão ser certificadas pelas autoridades educacionais competentes dos cursos frequentados.

§ 3º Para fins de cumulação dos casos de remição, as horas diárias de trabalho e de estudo serão definidas de forma a se compatibilizarem.

§ 4º O preso impossibilitado, por acidente, de prosseguir no trabalho ou nos estudos continuará a beneficiar-se com a remição.

§ 5º O tempo a remir em função das horas de estudo será acrescido de 1/3 (um terço) no caso de conclusão do ensino fundamental, médio ou superior durante o cumprimento da pena, desde que certificada pelo órgão competente do sistema de educação.

§ 6º O condenado que cumpre pena em regime aberto ou semiaberto e o que usufrui liberdade condicional poderão remir, pela frequência a curso de ensino regular ou de educação profissional, parte do tempo de execução da pena ou do período de prova, observado o disposto no inciso I do § 1º deste artigo.

§ 7º O disposto neste artigo aplica-se às hipóteses de prisão cautelar.

§ 8º A remição será declarada pelo juiz da execução, ouvidos o Ministério Público e a defesa. (NR)

\end{citacao}

Contudo, a remição além de funcionar como recompensa pela ordem e responsabilidade, incentivando através de atividades que contribuem para o retorno à vida em comunidade, tem como principal objetivo incentivar a reintegração do apenado à sociedade.\cite{lep}


